%----------------------------------------------------------------------------------------
%   USEFUL COMMANDS
%----------------------------------------------------------------------------------------

\newcommand{\dipartimento}{Dipartimento di Matematica ``Tullio Levi-Civita''}

%----------------------------------------------------------------------------------------
% 	USER DATA
%----------------------------------------------------------------------------------------

% Data di approvazione del piano da parte del tutor interno; nel formato GG Mese AAAA
% compilare inserendo al posto di GG 2 cifre per il giorno, e al posto di 
% AAAA 4 cifre per l'anno
\newcommand{\dataApprovazione}{Data}

% Dati dello Studente
\newcommand{\nomeStudente}{Luca}
\newcommand{\cognomeStudente}{Slongo}
\newcommand{\matricolaStudente}{2111009}
\newcommand{\emailStudente}{luca.slongo@studenti.unipd.it}
\newcommand{\telStudente}{+ 39 333 750 4455}

% Dati del Tutor Aziendale
\newcommand{\nomeTutorAziendale}{Aramis}
\newcommand{\cognomeTutorAziendale}{Todesco}
\newcommand{\emailTutorAziendale}{aramistodesco@etherealcodesrl.com}
\newcommand{\telTutorAziendale}{+ 39 334 159 1594}
\newcommand{\ruoloTutorAziendale}{Co-founder \& CTO}

% Dati dell'Azienda
\newcommand{\ragioneSocAzienda}{Ethereal Code SRL}
\newcommand{\indirizzoAzienda}{Via Gioacchino Rossini, San Zenone degli Ezzelini (TV)}
\newcommand{\sitoAzienda}{https://www.etherealcodesrl.com}
\newcommand{\emailAzienda}{info@etherealcodesrl.com}
\newcommand{\partitaIVAAzienda}{P.IVA 05585020265}

% Dati del Tutor Interno (Docente)
\newcommand{\titoloTutorInterno}{Prof.}
\newcommand{\nomeTutorInterno}{Marco}
\newcommand{\cognomeTutorInterno}{Zanella}

\newcommand{\prospettoSettimanale}{
Per ciascun argomento trattato, il tutor aziendale fornisce in via preliminare materiale di documentazione sui concetti di base, propedeutico alle attività pratiche della settimana.

\begin{itemize}

\item \textbf{Prima Settimana - Onboarding e Fondamenti (40 ore)}
\begin{itemize}
    \item Introduzione all'azienda, ai progetti e agli obiettivi del tirocinio;
    \item Configurazione editor e Git, presa in carico dell'ambiente di sviluppo Docker (Engine 29) pre-configurato;
    \item Ripasso JavaScript moderno (ES6+), asincronia, API;
    \item Introduzione al paradigma MVC e architettura web;
    \item Panoramica stack: Next.js 16, NestJS 11, PostgreSQL 18, Node.js 24 LTS;
\end{itemize}

\item \textbf{Seconda Settimana - Backend Base con NestJS (40 ore)}
\begin{itemize}
    \item Introduzione a NestJS 11 (controller, service, module);
    \item Creazione API REST base;
    \item Connessione a database PostgreSQL 18 (Supabase);
    \item CRUD semplice su entità del gestionale interno;
    \item Testing base delle API;
\end{itemize}

\item \textbf{Terza Settimana - Logica Applicativa e Auth (40 ore)}
\begin{itemize}
    \item Implementazione autenticazione JWT;
    \item Gestione utenti e ruoli base;
    \item Middleware, guard e validazioni;
    \item Introduzione a Redis 8 per caching;
\end{itemize}

\item \textbf{Quarta Settimana - Frontend con Next.js (40 ore)}
\begin{itemize}
    \item Introduzione a Next.js 16 (routing, componenti, fetch dati);
    \item Stilizzazione delle interfacce con Tailwind CSS 4;
    \item Creazione interfaccia per il gestionale interno;
    \item Integrazione frontend-backend tramite API;
    \item Gestione stato e autenticazione lato client;
    \item Introduzione a Vercel come piattaforma di deploy del frontend, con preview environments;
\end{itemize}

\item \textbf{Quinta Settimana - Introduzione SaaS e Multitenancy (40 ore)}
\begin{itemize}
    \item Concetti base di SaaS e architettura multitenant;
    \item Strutturazione backend per gestione tenant;
    \item Separazione dati e logica multi-cliente;
    \item Estensione del progetto esistente;
\end{itemize}

\item \textbf{Sesta Settimana - Infrastruttura e Container (40 ore)}
\begin{itemize}
    \item Comprensione dell'ambiente Docker Engine 29 pre-configurato (compose, networking);
    \item Introduzione a Kubernetes 1.36 nelle distribuzioni k3s e k8s tramite l'ambiente di produzione AWS pre-impostato;
    \item Comprensione del flusso di deploy dei servizi backend su AWS e dei frontend su Vercel;
    \item Utilizzo S3 per gestione file;
\end{itemize}

\item \textbf{Settima Settimana - Architetture Scalabili (40 ore)}
\begin{itemize}
    \item Introduzione a RabbitMQ 4 e code di messaggi;
    \item Gestione processi asincroni;
    \item API Gateway con Kong Gateway 3.10 LTS e concetti di microservizi;
    \item Ottimizzazione performance e caching;
\end{itemize}

\item \textbf{Ottava Settimana - Consolidamento e Conclusione (40 ore)}
\begin{itemize}
    \item Revisione completa dei progetti sviluppati;
    \item Debugging e miglioramenti;
    \item Documentazione tecnica;
    \item Presentazione finale del lavoro svolto;
\end{itemize}

\end{itemize}
}

% Indicare il totale complessivo (deve essere compreso tra le 300 e le 320 ore)
\newcommand{\totaleOre}{320}

\newcommand{\obiettiviObbligatori}{
    \item \underline{\textit{O01}}: Comprendere le basi dello sviluppo web moderno lato backend utilizzando Node.js 24 LTS e NestJS 11;

    \item \underline{\textit{O02}}: Sviluppare e testare API REST complete con operazioni CRUD integrate a database PostgreSQL 18 (Supabase);

    \item \underline{\textit{O03}}: Implementare un sistema di autenticazione basato su JWT con gestione utenti e autorizzazioni di base;

    \item \underline{\textit{O04}}: Realizzare una semplice interfaccia frontend con Next.js 16 e Tailwind CSS 4 integrata alle API sviluppate;

    \item \underline{\textit{O05}}: Utilizzare l'ambiente Docker Engine 29 pre-configurato per l'esecuzione e la gestione dell'applicazione in sviluppo;
}

\newcommand{\obiettiviDesiderabili}{
    \item \underline{\textit{D01}}: Comprendere i concetti di architettura SaaS e multitenancy;

    \item \underline{\textit{D02}}: Introdurre meccanismi di caching tramite Redis 8 per migliorare le performance;

    \item \underline{\textit{D03}}: Acquisire familiarità con sistemi di messaggistica asincrona (RabbitMQ 4);

    \item \underline{\textit{D04}}: Comprendere il ruolo di un API Gateway (Kong Gateway 3.10 LTS) in un'architettura a microservizi;
}

\newcommand{\obiettiviFacoltativi}{
    \item \underline{\textit{F01}}: Comprendere il flusso di deploy sull'infrastruttura cloud pre-configurata, con servizi backend su AWS e frontend su Vercel;

    \item \underline{\textit{F02}}: Comprendere i concetti fondamentali di orchestrazione con Kubernetes 1.36 nelle distribuzioni k3s e k8s tramite l'ambiente di produzione pre-impostato;
    
    \item \underline{\textit{F03}}: Contribuire alla documentazione tecnica del progetto;
}